My CR5 deployment and Hermite action-tokenization research have taught me to analyze execution semantics and representation fidelity, but they also expose a capability gap: I still need a disciplined way to unite robot learning, control, and physical experimentation in one validation loop. CUHK's MSc in Robotics addresses that gap directly. If offered in my intake, ROSE5760 Robot Learning would help me examine how policies acquire behavior, while ENGG5402 Advanced Robotics would let me relate those policies to the control structures that execute them. This pairing matches the problem I encountered when non-blocking CR5 commands used intermediate arm states and generated unstable corrections.

I would then use ROSE5720 Experimental Robotics to test synchronization and short-horizon execution on hardware rather than infer execution reliability from offline metrics. Subject to project approval and availability, ROSE5910 M.Sc. Project would give me a setting to investigate two independent questions without conflating them: whether Hermite tokenization preserves trajectory structure, and whether CR5 execution preserves state-action consistency. I would contribute an integration-first perspective grounded in real data and control interfaces, and develop toward robotics research and development work on learned policies that remain stable when transferred from representation to a physical robot.
